\section{Large deviations and concentration}
\label{sec: concentration inequalities}

In the previous section was concerned with the distribution of \emph{typical deviations} of
the mean, which were of order \(\frac1{\sqrt{n}}\). The purpose of this section is two-fold. First, we want to understand
how fast the probability of \emph{large deviations} of constant order \(1\) decay with \(n\).
Second, we want more quantitative bounds on the deviation probability.
Observe that the central limit theorem only gives an asymptotic statement:
``If \(n\to\infty\), then...''. However in practice \(n\) is always finite.
And we can never be sure that the normal distribution is really a good
approximation of the distribution at finite \(n\).
\begin{quotation}
    ``In the long run we are all dead'' -- John Maynard Keynes 
\end{quotation}
It is a lucky coincidence that the tools to analyze large deviations are also
better suited to provide quantitative bounds. Quantitative versions of the CLT
also exist, but they come at the price of much greater difficulty.

With a quantitative bound we mean that we want an explicit upper bound
\[
    \prob[\big]{\norm{X_n - X} \ge \epsilon} \le \delta
\]
for a fixed \(n\) instead of just the statement that this probability goes to zero as
\(n\to\infty\). Such statements are also know as ``probably, approximately
correct'' (PAC) statements, that is
\begin{quote}
``with probability at least \(1-\delta\) (probably), we have \(\norm{X_n - X} <
\epsilon\) (approximately) for a specific \(n\in \nat\)''.
\end{quote}

\subsection{A real valued introduction}

Recall that the only tool we have used so far to bound probabilities is the
Markov inequality (Theorem \ref{thm: markov inequality})
\[
    \prob{X \ge x} \le \tfrac{\E{h(X)}}{h(x)}.
\]
The choice \(h(x) = e^{tx}\), with \(t\ge 0\) such that \(h\) is
non-decreasing, results in the ``\textbf{Chernoff bound}''
\[
    \prob{X \ge x}
    \le \frac{\E{e^{t  X}}}{e^{t x}}
    = \frac{\MGF_X(t)}{e^{t x}}
    = \exp\bigl(-\bigl(t x - \underbrace{\log \MGF_X(t)}_{=\CGF_X(t)}\bigr)\bigr).
\]
If the moment generating function of \(X\) is not finite, then this bound is
of course useless. Note that we cannot use the characteristic function here,
since the Markov inequality requires non-negative, and in particular
real-valued, functions \(h\). The function \(\CGF_X\) is the ``\textbf{cumulant
generating function}'' (CGF) or ``log-laplace transform'' that will be our main object
of study in this section.

The key feature of the Chernoff bound is that it provides us with a parameter \(t\)
that is not inherent to the problem of bounding \(\prob{X \ge a}\). This allows for
an optimization over \(t\), resulting in the \textbf{Legendre transform}
\[
    \CGF^+_X(x) \coloneq \sup_{t\ge 0} \bigl( t x - \CGF_X(t) \bigr).
\]
The optimized bound is then given by
\begin{equation}
    \label{eq: chernoff bound}    
    \prob{X \ge x} \le \exp\bigl(-\CGF^+_X(x)\bigr).
\end{equation}

\begin{example}[The normal distribution]
    \label{ex: chernoff bound for normal distribution}
    Let \(X \sim \normal{\mu, \sigma^2}\). 
    Recall that the MGF of the normal distribution was
    \(\MGF_X(t) = \exp(\mu t + \frac12 \sigma^2 t^2)\) (Definition \ref{def:
    univariate normal distribution}). Its CGF is then given by the quadratic
    polynomial
    \[
        \CGF_X(t) = \mu t + \tfrac12 \sigma^2 t^2.
    \]
    Thus the Legendre transform is
    \[\begin{aligned}
        \CGF_X^+(x)
        = \sup_{t\ge 0} \bigl(t x - \mu t - \tfrac12 \sigma^2 t^2\bigr)
        &= \sup_{t\ge 0} \bigl(t (x - \mu) - \tfrac12 \sigma^2 t^2\bigr)
        \\
        &= \begin{cases}
        \frac{(x - \mu)^2}{2\sigma^2}
        \qquad &\text{for } x \ge \mu,
        \\
        0 &\text{for } x \le \mu.
        \end{cases}
    \end{aligned}\]
    The last step follows from simple optimization. The first order condition
    results in the maximizer \(t^\star = \frac{x - \mu}{\sigma^2}\), which is greater
    zero if and only if \(x \ge \mu\).
    Observe that if \(\CGF_X^+(x) = 0\), then the Chernoff bound \eqref{eq: chernoff bound}
    is \(1\) which is not very helpful since probabilities are always bounded by \(1\).
    The only useful case is thus \(x \ge \mu\), in which we get that
    the tail of the normal distribution decays like a squared exponential in the
    distance \((x-\mu)\) to the mean, that is
    \[
        \prob{X \ge x}
        \le \exp\bigl(- \tfrac{(x - \mu)^2}{2\sigma^2}\bigr)
        \qquad\text{for } x \ge \mu.
    \]
    Or equivalently, with \(\epsilon = x - \mu > 0\) we get a quickly decaying bound in \(\epsilon\)
    \[
        \prob{X - \mu \ge \epsilon}
        \le \exp(-\tfrac{\epsilon^2}{2\sigma^2})
        \qquad\text{for } \epsilon \ge 0.
    \]
\end{example}

\paragraph*{Towards multivariate bounds}
Since \(\CGF_X^+\) is defined as a supremum over all \(t\ge 0\), it does not
translate well into multiple dimension where the ordering is lost.
We have seen in the example above that the maximizer \(t^\star\) is positive
whenever \(x\ge \mu\) or equivalently \(\epsilon = x - \mu \ge 0\).
This implies that for \(x\ge \mu\) the function \(\CGF_X^+(x)\) coincides
with the \textbf{convex conjugate}, also known as ``Fenchel-Legendre transform''
of \(\CGF_X\) given by
\[
    \CGF_X^*(x) \coloneq \sup_{t\in \real} \bigl(t x - \CGF_X(t)\bigr)
    \qquad \Bigl(\ge \CGF_X^+(x)= \sup_{t \ge 0} \bigl(t x - \CGF_X(t)\bigr)\Bigr).
\]
\begin{warning}
    \label{warn: difference between CGF star and plus}
    For \(x\le \mu\) the bound \(\prob{X \ge x} \le \exp(-\CGF_X^*(x))\) is
    false! Here \(\CGF_X^*(x) > \CGF_X^+(x) = 0\).  However since
   \(\CGF_X^+(x)=0\) is a useless bound anyway, we simply need to avoid this
    case.
\end{warning}

% To understand the convex conjugate for \(a\le \mu\), observe that in this case
% \(-a \ge -\mu = \E{-X}\). Consequently,
% \[
% \begin{aligned}
%     \CGF_{-X}^+(-a)
%     = \sup_{t\ge 0} \bigl(t (-a) - \CGF_{-X}(t)\bigr)
%     &= \sup_{t\ge 0} \bigl(-ta - \CGF_{X}(-t)\bigr)
%     \\
%     &= \sup_{s\le 0} \bigl(sa - \CGF_{X}(s)\bigr)
%     = \CGF_X^*(a),
% \end{aligned}
% \]
% where we use \(\CGF_{-X}(t) = \log\E{e^{t(-X)}}= \CGF_X(-t)\).
% So in essence
% \begin{alignat}{3}
%     \prob{X \ge a} &\le \exp\bigl(-\CGF_X^*(a)\bigr) &\text{for }a &\ge \mu,
%     \\ 
%     \prob{X \le a} = \prob{-X \ge -a} &\le \exp\bigl(-\CGF_{X}^*(a)\bigr)\qquad &\text{for } a &\le \mu.
% \end{alignat}
% Why is this useful? For \(a \le \mu \le b\) consider
% \[
%     \prob{X \in (a,b)^\complement}
%     = \prob{X \le a} + \prob{X \ge b}
% \]

% For concentration bounds we will simply use
% \[\begin{aligned}
%     \prob{\abs{X} \ge a}
%     &= \prob{X \ge a} + \prob{-X \ge a}
%     \\
%     &\le \exp\bigl(-\CGF^+_X(a)\bigr) + \exp\bigl(-\CGF^+_{-X}(a)\bigr)
% \end{aligned}\]

\subsection{Large deviations}
\marginpar{\red{Lecture 6}}

We proceed with a more formal introduction of the cumulant generating function
and its properties in multiple dimensions.

\begin{definition}[Cumulant generating function]
    The \textbf{cumulant generating function} (CGF) of a random variable \(X\)
    in a Hilbertspace \(\hilbert\) is defined as the logarithm of its moment generating
    function
    \[\tag{CGF}
        \CGF_X(t) \coloneq \log \MGF_X(t)
        = \log \E{e^{\scp{t, X}}}.
    \]
    Due to \(\exp(x) > 0\) the MGF \(\MGF_X\) is always defined with values in \((0,
    \infty]\). Thus the CGF is defined for all \(t\in \hilbert\) with values in \((-\infty,
    \infty]\). We denote its effective domain
    by
    \[
        \tag{effective domain}
        \cD_{\CGF_X} \coloneq \set{t \in \hilbert: \CGF_X(t) < \infty} = \cD_{\MGF_X} \supseteq \set{0}.
    \]
    We further define the functions for \(t,a\in \hilbert\)
    \begin{align*}
        \tag{concave-convex helper}
        \Phi_X(t, a) &\coloneq \scp{t, a} - \CGF_X(t)
        \\
        \tag{convex conjugate of \(\CGF_X\)}
        \CGF_X^*(a) &\coloneq \sup_{t\in \hilbert} \Phi(t, a)
    \end{align*}
\end{definition}

\begin{lemma}[Properties of the CGF]
    \label{lem: properties of CGF}
    Let \(X\) be a random variable in a Hilbertspace \(\hilbert\). Then
    \begin{enumerate}[label=\normalfont{(\arabic*)}, noitemsep]
        \item\label{it: CGF(0)=0} \(\CGF_X(0) = 0\),
        \item\label{it: CGF convex} \(\CGF_X\) is convex,
        \item\label{it: CGF lower semicontinuous} \(\CGF_X\) is lower semicontinuous, i.e.\ \(\liminf\limits_{n\to\infty} \CGF_X(t_n) \ge \CGF_X(t)\) for any \(t_n\to t\),
        in other words: the CGF takes the bottom value at jump points

        \item\label{it: CGF effective domain is convex} \(\cD_{\CGF_X}\) is a convex set that contains the origin \(0\),
    \end{enumerate} 
    The CGF interacts nicely with sums (the key in the analysis of the mean)
    \begin{enumerate}[label=\normalfont{(\arabic*)}, noitemsep, resume]
        \item\label{it: CGF sum independent} \(\CGF_{X+Y}(t) = \CGF_X(t) + \CGF_Y(t)\) for \(X,Y\) independent and all \(t\in \hilbert\).
        \item\label{it: CGF affine} \(\CGF_{aX + b}(t) = \scp{t, b} + \CGF_X(a t)\)
        for any \(a\in \real\), \(b, t\in \hilbert\).
    \end{enumerate}
    Let \(\hilbert=\real^\dims\) and assume \(U \subseteq \cD_{\CGF_X}\) is an open
    set with \(0\in U\), then
    \begin{enumerate}[label=\normalfont{(\arabic*)}, noitemsep, resume]
        \item\label{it: CGF analytic} \(\CGF_X\) is analytic (a power series) on \(U\), specifically for \(\dims=1\)
        \[
            \CGF_X(t)
            = \sum_{k=1}^\infty \CGF_X^{(k)}(0)\frac{t^k}{k!}
            \qquad \forall t\in O.
        \]
        The coefficient \(\cumulant_k(X) \coloneq \CGF_X^{(k)}(0)\) is called the \(k\)-th \textbf{cumulant} of \(X\).
        For \(X,Y\) independent we have by \ref{it: CGF sum independent},
        \[
            \cumulant_k(X+Y) = \cumulant_k(X) + \cumulant_k(Y).
        \]
        \item\label{it: CGF gradient and Hessian} \(\nabla \CGF_X(0) = \E{X}\) and \(\nabla^2\CGF_X(0) = \Cov{X}\).
    \end{enumerate}
\end{lemma}
\begin{proof}
    \ref{it: CGF(0)=0}, \ref{it: CGF sum independent} and \ref{it: CGF affine}
    are easy to verify and also follow from the equivalent properties of the MGF
    (Exercise \ref{ex: properties moment generating functions}). For \ref{it:
    CGF analytic} recall that the MGF is analytic on the interior of its
    effective domain (Lemma \ref{lem: moment generating function representation}
    and Exercise \ref{ex: multivariate
    moment generating function}) and the \(\log\) is analytic on \((0,
    \infty)\), therefore the CGF is also analytic there. Then \ref{it: CGF
    gradient and Hessian} follows from \(\log \MGF_X(t)\) and the chain rule.
    
    Convexity of \(\CGF_X\), \ref{it: CGF convex}, follows from Hölder's inequality (Exercise \ref{ex: hölder's inequality}).
    Specifically, for any \(t, s\in \hilbert\) and any \(\alpha\in[0,1]\) we choose
    \(p = \frac1\alpha\) and \(q = \frac1{1-\alpha}\) such that \(\frac1p
    + \frac1q = 1\), then
    \[\begin{aligned}[b]
        \CGF_X(\alpha t + (1-\alpha) s)
        &= \log \E{e^{\scp{\alpha t + (1-\alpha) s, X}}}
        \\
        \overset{\text{Hölder}}&\le \log \E{(e^{ \alpha\scp{t, X}})^{p}}^{\frac1p} \E{(e^{(1-\alpha)\scp{s, X}})^{q}}^{\frac1q}
        \\
        \overset{p=\frac1\alpha}&= \log \E{e^{\scp{t, X}}}^{\alpha} \E{e^{\scp{s, X}}}^{1-\alpha}
        = \alpha \CGF_X(t) + (1-\alpha) \CGF_X(s).
    \end{aligned}
    \]
    The convexity of the effective domain, \ref{it: CGF effective domain is convex}, follows directly.
    Lower semicontinuity of \(\CGF_X\), \ref{it: CGF lower semicontinuous}, follows from Fatou's lemma (measure
    theory):
    \[\begin{aligned}[b]
        \liminf_{n\to \infty} \CGF_X(t_n)
        &= \log \liminf_{n\to \infty} \E{\exp(\scp{t_n, X})}
        \\
        \overset{\text{Fatou}}&\ge \log \E{\liminf_{n\to \infty} \exp(\scp{t_n, X})}
        = \CGF_X(t).
    \end{aligned}
    \qedhere
    \]
\end{proof}

\begin{remark}[Often the CGF is non-negative]
    If \(\E{X} = 0\), then the slope of the CGF at \(0\) is zero, that is \(\nabla \CGF_X(0) = 0\).
    Since the CGF is convex and the second derivative \(\Cov{X}\) is positive definite,
    this implies it must be non-negative everywhere. For \(\E{X} \neq 0\)
    this is obviously not the case.
\end{remark}

Observe that \(\Phi_X\) is thus a concave-convex function since \(\CGF_X\) is convex.
This explains its name and this property is crucial for the following bound.

\begin{mdframed}[innertopmargin=0pt]
\begin{proposition}
    \label{prop: general LDP upper}
    Let \(A\subseteq \hilbert\) be a \underline{compact convex} set, then 
    \[
        \prob{X \in A} \le \exp\bigl(-\inf_{a\in A} \CGF_X^*(a)\bigr).
    \]
\end{proposition}
\end{mdframed}

\begin{proof}
    Then
    \(
        \scp{t, x} - \inf_{a\in A} \scp{t, a} \ge 0
    \)
    for all \(x\in A\) and any \(t\in \hilbert\). Since \(e^\lambda \ge 1\)
    for all \(\lambda \ge 0\) we have
    \[\begin{aligned}
        \prob{X \in A} = \E{\ind{X \in A}}
        &\le \E[\Big]{\exp\bigl(\scp{t, X} - \inf_{a\in A} \scp{t, a}\bigr)}
        \\
        &= \MGF_X(t)\exp\bigl(- \inf_{a\in A} \scp{t, a}\bigr) 
        \\
        &= \exp\Bigl(- \bigl(\inf_{a\in A} \scp{t, a}-\CGF_X(t)\bigr)\Bigr) 
        = \exp(- \inf_{a\in A} \Phi_X(t, a)).
    \end{aligned}
    \]
    Since this bound holds for any \(t\in \hilbert\) we have by Sion's
    minimax theorem (Fact \ref{fact: sion's minimax theorem}) and Lemma
    \ref{lem: check sion's conditions}
    \[\begin{aligned}[b]
        \prob{X \in A}
        &\le \exp\Bigl(- \sup_{t \in \hilbert} \inf_{a\in A} \Phi_X(t, a)\Bigr)
        \\
        \overset{\text{Sion}}&=
        \exp\Bigl(- \inf_{a\in A} \sup_{t \in \hilbert} \Phi_X(t, a)\Bigr)
        = \exp\bigl(- \inf_{a\in A} \CGF_X^*(a)\bigr).
    \end{aligned}
    \qedhere
    \]
\end{proof}

\begin{fact}[Sion's minimax theorem {\citep[e.g.][]{komiyaElementaryProofSions1988}}]
    \label{fact: sion's minimax theorem}
    Let \(\cX\) be a \underline{compact convex} subset of a topological vectorspace and
    \(\cY\) a convex subset of a topological vectorspace. Let
    \(f \colon \cX \times \cY \to \real\) be a function such that
    \begin{enumerate}[label=(\roman*), noitemsep]
        \item for each \(y \in \cY\), \(f(\cdot, y)\) is upper semicontinuous and quasi-concave on \(\cX\),
        \item for each \(x \in \cX\), \(f(x, \cdot)\) is lower semicontinuous and quasi-convex on \(\cY\).
    \end{enumerate}
    Then
    \[
        \sup_{x \in \cX} \inf_{y \in \cY} f(x, y)
        = \inf_{y \in \cY} \sup_{x \in \cX} f(x, y);
    \]
\end{fact}

\begin{lemma}[Check Sion's conditions]
    \label{lem: check sion's conditions}
    \(\Phi(t, a) = \scp{t, a} - \CGF_X(t)\) satisfies the conditions of Sion's minimax theorem.
\end{lemma}
\begin{proof}
    The function \(\Phi_X\) is trivially continuous and convex in \(a\) since
    the inner product is continuous and linear in \(a\).
    Since \(\CGF_X\) is lower semicontinuous and convex (Lemma \ref{lem:
    properties of CGF}) it follows that
    the function \(\Phi_X\) is upper semicontinuous and concave in \(t\).
\end{proof}

In finite dimensions we can extend Proposition \ref{prop: general LDP upper} to
\emph{closed} convex sets.

\begin{mdframed}[innertopmargin=0pt]
\begin{theorem}[General Chernoff bound]
    \label{thm: chernoff bound in R^d}
    Let \(A\subseteq \real^\dims\) be a \underline{closed}, \underline{convex} set, then 
    \[
        \prob{X \in A} \le \exp\bigl(-\inf_{a\in A} \CGF_X^*(a)\bigr).
    \]
\end{theorem}
\end{mdframed}
\begin{proof}
    Let \((K_n)_{n\in \nat}\) be an increasing sequence of convex compact sets
    with \(K_n \uparrow \real^\dims\) (this does not necessarily exist in infinite
    dimension!). Then by continuity of probability measures
    (Remark \ref{rem: continuity of probability})
    \[\begin{aligned}[b]
        \prob{X \in A}
        &\le \lim_{n\to \infty} \prob{X \in A \cap K_n} + \underbrace{\prob{X \notin K_n}}_{\to 0}
        \\
        &\le \lim_{n\to \infty} \exp\bigl(- \inf_{a\in A \cap K_n} \CGF_X^*(a)\bigr)
        \le \exp\bigl(- \inf_{a\in A} \CGF_X^*(a)\bigr).
    \end{aligned}
    \qedhere
    \]
\end{proof}

\paragraph*{Comparison with the one-dimensional motivation}
Let us compare the generalized Chernoff bound with the Chernoff bound
we obtained in the one-dimensional case. The one dimensional bound \eqref{eq: chernoff bound}
was
\[
    \prob{X \ge x} \le \exp\bigl(-\CGF_X^+(x)\bigr)
    \Bigl(= \exp\bigl(-\CGF_X^*(x)\bigr) \qquad \text{for } x \ge \E{X}\Bigr).
\]
Using the closed, convex set \(A=[x, \infty)\) in the generalized Chernoff bound we get
\[
    \prob{X \ge x} \le \exp\bigl(-\inf_{a\ge x} \CGF_X^*(a)\bigr).
\]
Since \(\CGF_X^*(x) \ge \inf_{a\ge x} \CGF_X^*(a)\) it may appear
that we obtained a worse bound. However the key to understand here is that the
function \(\CGF_X^*(a)\) is monotonically increasing after \(\E{X}\) (Lemma \ref{lem: properties of cgf*}) 
and therefore
\[
    \CGF_X^*(x) = \inf_{a\ge x} \CGF_X^*(a) \quad \text{for all } x \ge \E{X}.
\]
More specifically, we will see in the following that the expected value \(\E{X}\)
is the minimizer of \(\CGF_X^*(a)\) and that \(\CGF_X^*(a)\) is also convex and therefore
monotonically increasing for \(a \ge \E{X}\). Observe that for \(x< \E{X}\)
the minimizer is the expected value and we do not have the problem to ensure \(x
\ge \E{X}\) anymore (cf.~Warning \ref{warn: difference between CGF star and
plus}).

\paragraph*{Understanding the convex conjugate} Recall that a convex function
\(\Lambda\) has the property that it lies above its supporting hyperplanes. That is 
\[
    \Lambda(t) \ge \scp[\big]{t-s, \nabla \Lambda(s)} + \Lambda(s)
    = \underbrace{\scp[\big]{t, \nabla \Lambda(s)}}_{\text{hyperplane}}
    - \underbrace{\bigl(\scp[\big]{s, \nabla \Lambda(s)} - \Lambda(s)\bigr)}_{\eqcolon c(s) \quad\text{(intercept)}}.
\]
The convex conjugate \(\Lambda^*\) is essentially the map from the gradient to
the intercept, \(\nabla \Lambda(s) \mapsto c(s)\). That is
\[
    \Bigl(\Lambda(t) \ge \scp{t, a} - c\qquad \forall t \in \hilbert\Bigr)
    \iff c \ge \bigl(\sup_{t \in \hilbert} \scp{t, a} - \Lambda(t)\bigr) = \Lambda^*(a).
\]
So the smallest intercept that ensures the hyperplane is below the function is
exactly the convex conjugate. Being the smallest intercept it ensure the hyperplane
is touching the function at some point (cf.~\(t=s\) above). In particular,
if \(\Lambda\) is differentiable, then
\begin{equation}
    \label{eq: gradient of cgf star}    
    \Lambda^*\bigl(\nabla\Lambda(t)\bigr)
    = c(t)
    = \scp[\big]{t, \nabla \Lambda(t)} - \Lambda(t).
\end{equation}
Recall that if the essential domain of \(\CGF_X\) contains an open
neighborhood of \(0\), then \(\nabla \CGF_X(0) = \E{X}\) and consequently
\begin{equation}
    \label{eq: cgf star at expected value}    
    \CGF_X^*\bigl(\E{X}\bigr) = \CGF_X^*\bigl(\nabla \CGF_X(0)\bigr) = \scp{0, \nabla \CGF_X(0)} - \CGF_X(0) = 0.
\end{equation}
So up to the issue with the domain it is sufficient to prove that \(\CGF_X\)
is non-negative and convex to get that \(\E{X}\) is its minimizer.


\begin{lemma}[Properties of the convex conjugate]
    \label{lem: properties of cgf*}
    The convex conjugate \(\Lambda^*\) of an arbitrary function \(\Lambda\)
    satisfies
    \begin{enumerate}[label=\normalfont{(\arabic*)}, noitemsep]
        \item\label{it: cgf* convex} \(\Lambda^*\) is convex,
        \item\label{it: cgf* lower semicontinuous} \(\Lambda^*\) is lower semicontinuous,
        \item\label{it: cgf* duality} \(\Lambda^{**} = \Lambda\) if \(\Lambda\) is convex and lower semicontinuous (Fenchel-Moreau Thm.),
        \item\label{it: cgf* gradient} If \(\Lambda\) is
        differentiable in \(t\in \hilbert\), then \(\Lambda^*(\nabla \Lambda(t)) = \scp{t, \nabla \Lambda(t)} - \Lambda(t)\)
        \item\label{it: cgf* minimum at expected value} If \(\Lambda(0) = 0\), then \(\Lambda^*(a) \ge 0\) for all \(a\in \hilbert\).
    \end{enumerate}
    Moreover for \(\Lambda = \CGF_X\) we have
    \begin{enumerate}[label=\normalfont{(\arabic*)}, noitemsep, resume]
        \item\label{it: cgf* trivial essential domain} If \(\cD_{\CGF_X} = \set{0}\), then \(\CGF_X \equiv 0\).
        \item\label{it: cgf* open essential domain} If \(U \subseteq \cD_{\CGF_X}\) is an open set with \(0\in U\), then
        \(\CGF_X^*(\E{X})=0\).
        In particular, the expectation \(\E{X}\) is a minimizer of \(\CGF_X^*\).
    \end{enumerate}
\end{lemma}

\begin{remark}
    Between \ref{it: cgf* trivial essential domain} and \ref{it: cgf* open
    essential domain} there is a gap that is very annoying to treat formally,
    but offers little insight.
    See \fxwarning{reference or exercise?}
\end{remark}

\begin{proof}
    For \ref{it: cgf* convex} choose \(\lambda\in [0,1]\) and \(a, b\in \hilbert\). Then
    \[\begin{aligned}
        \Lambda^*(\lambda a + (1-\lambda) b)
        &= \sup_{t\in \hilbert} \bigl(\lambda \scp[\big]{t, a} + (1-\lambda) \scp[\big]{t, b} - \Lambda(t)\bigr)
        \\
        &\le \lambda \sup_{t\in \hilbert} \bigl(\scp{t, a} - \Lambda(t)\bigr)
        + (1-\lambda) \sup_{t\in \hilbert} \bigl(\scp{t, b} - \Lambda(t)\bigr)
        \\
        &= \lambda \Lambda^*(a) + (1-\lambda) \Lambda^*(b).
    \end{aligned}
    \]
    For \ref{it: cgf* lower semicontinuous} let \(a_n \to a\). For any fixed \(t\in \hilbert\) we have
    \[
        \liminf_{n\to \infty} \Lambda^*(a_n)
        \overset{\text{Def.~}\Lambda^*}\ge \liminf_{n\to \infty}\; \scp{t, a_n} - \Lambda(t)
        = \scp{t, a} - \Lambda(t).
    \]
    Since this holds for any \(t\in \hilbert\), it also holds for the supremum
    \[
        \liminf_{n\to \infty} \Lambda^*(a_n) \ge \sup_{t\in \hilbert} \bigl(\scp{t, a} - \Lambda(t)\bigr) = \Lambda^*(a).
    \]
    We will not prove \ref{it: cgf* duality}. See \fxwarning{reference!}. We have already proven
    \ref{it: cgf* gradient} in \eqref{eq: gradient of cgf star}. For \ref{it: cgf* minimum at expected value} observe that
    \[
        \Lambda^*(a) = \sup_{t\in \hilbert} \bigl(\scp{t, a} - \Lambda(t)\bigr)
        \ge \scp{0, a} - \Lambda(0) = 0.
    \]
    \ref{it: cgf* trivial essential domain} follows directly from the definition
    as \(\scp{t, a} - \CGF_X(t) = -\infty\) for all \(t \neq 0\). Consequently
    the supremum chooses \(t=0\).  \ref{it: cgf* open essential domain} was already
    shown in \eqref{eq: cgf star at expected value}.
\end{proof}

\paragraph*{Large deviation principle for the mean}

Since they are closely related to the MGF, the CGF also interacts very well with
independent sums.

\begin{lemma}
    Let \(X_1, \ldots, X_n\) be independent copies of the random variable \(X\)
    and \(\mean{X}_n\) their mean, then
    \begin{equation}
        \label{eq: cumulant generating function of the mean}
        \CGF_{\mean{X}_n}(t) = n \CGF_X\bigl(\tfrac{t}n\bigr)
        \qquad\text{and}\qquad
        \CGF_{\mean{X}_n}^*(t) = n \CGF_X^*(t).
    \end{equation}
\end{lemma}
\begin{proof}
    The first claim follows by simple calculation and the nice interaction of
    the exponential with independent sums
    \[\begin{aligned}
        \CGF_{\mean{X}_n}(t)
        &= \log \E[\Big]{\exp\Bigl(\scp[\Big]{t, \frac1n \sum_{i=1}^n X_i}\Bigr)}
        \\
        &= \log \E[\Big]{\prod_{i=1}^n e^{\scp{\frac{t}n, X_i}}}
        \overset{\iid}= \log \MGF_X\bigl(\tfrac{t}n\bigr)^n
        = n \CGF_X\bigl(\tfrac{t}n\bigr).
    \end{aligned}
    \]
    The second claim follows from the first by a change of variable
    \[
        \CGF_{\mean{X}_n}^*(a)
        = \sup_{t\in \hilbert} \bigl(n \scp{\tfrac{t}n, a} - n \CGF_X\bigl(\tfrac{t}n\bigr)\bigr)
        \overset{s=\frac{t}n}= n \sup_{s\in \hilbert} \bigl(\scp{s, a} - \CGF_X(s)\bigr)
        = n \CGF_X^*(a).
        \qedhere
    \]
\end{proof}

The Lemma above immediately implies exponential convergence to the mean
\begin{mdframed}[innertopmargin=0pt]
\begin{corollary}
    Let \(A\subseteq \real^\dims\) be a closed convex set, then
    \[
        \prob[\big]{\mean{X}_n \in A} \le \exp\bigl(-n \inf_{a\in A} \CGF_X^*(a)\bigr).
    \]
\end{corollary}
\end{mdframed}
In the one-dimensional case we get concentration bounds from this:
\begin{mdframed}[innertopmargin=0pt]
\begin{theorem}[Concentration bound for the mean]
    Let \(X_1, \ldots, X_n\) be independent copies of the random variable \(X\)
    with expectation \(\mu=\E{X}\) in \(\real\) and \(\mean{X}_n\) their mean,
    then for any \(t > 0\)
    \[
        \prob[\Big]{\abs[\big]{\mean{X}_n - \mu} \ge t} \le 2 \exp\Bigl(-n \inf_{\abs{a-\mu}\ge t} \CGF_X^*(a)\Bigr).
    \]
\end{theorem}
\end{mdframed}

\begin{proof}
    Since the sets \((-\infty, \mu - t]\) and \([\mu + t, \infty)\) are closed
    and convex we have
    \[\begin{aligned}[b]
        \prob[\big]{\abs{\mean{X}_n - \mu} \ge t}
        &= \prob{\mean{X}_n \ge t + \mu} + \prob{\mean{X}_n \le \mu - t}
        \\
        &\le \exp\bigl(-n \inf_{a \ge t+\mu}\CGF_X^*(a)\bigr) + \exp\bigl(-n \inf_{a \le \mu - t}\CGF_{X}^*(a)\bigr)
        \\
        &\le 2 \exp\Bigl(-n \min\set[\Big]{\inf_{a-\mu \ge t}\CGF_X^*(a), \inf_{a-\mu \le -t}\CGF_X^*(a)}\Bigr)
        \\
        &= 2 \exp\Bigl(-n \inf_{\abs{a-\mu}\ge t} \CGF_X^*(a)\Bigr).
    \end{aligned}
    \qedhere
    \]
\end{proof}

Observe that the factor \(2\) appears since the set \((-\infty, \mu - t]
\cup [\mu + t, \infty)\) is not convex and we had to split it into two convex
parts. This has however no bearing on the exponential rate of convergence:
\[
    \frac1n \log \prob[\big]{\abs{\mean{X}_n - \mu} \ge t}
    \le \underbrace{\frac{\log{2}}n}_{\to 0} - \inf_{\abs{a-\mu}\ge t} \CGF_X^*(a)
\]
With this less `sharp' view, we can thereby remove the convexity assumption on the set.
In particular this exponential rate is tight.

\begin{mdframed}[innertopmargin=0pt]
\begin{fact}[Cramér]
    Let \(X_1, \ldots, X_n\) be independent copies of the random variable \(X\) 
    in \(\real^\dims\) and \(\mean{X}_n\) their mean. Let \(A\) be a closed (but
    not necessarily convex) set, then
    \[
        \limsup_{n\to\infty} \frac1n \log \prob[\big]{\mean{X}_n \in A} \le -\inf_{a\in A} \CGF_X^*(a).
    \]
    Moreover for any open set \(U\)
    \[
        \liminf_{n\to\infty} \frac1n \log \prob[\big]{\mean{X}_n \in U} \ge -\inf_{a\in U} \CGF_X^*(a).
    \]
    In particular, for any set \(B\) that does not have isolated points, i.e.~\(\closure{\interior{B}} = \closure{B}\),
    \[
        \lim_{n\to\infty} \frac1n \log \prob[\big]{\mean{X}_n \in B} = -\inf_{a\in B} \CGF_X^*(a).
    \]
\end{fact}
\end{mdframed}

The lower bound shows that the exponential rate is tight. Since \(U\)
is open it is proven by taking a small ball around the point \(a\in U\) that minimizes \(\CGF_X^*(a)\) and
lower bounding the probability of this ball.
The upper bound is obtained for
arbitrary closed sets by covering the set with convex sets and using the fact
that the factor (like the \(2\) we got) goes to zero when the logarithm and
\(\frac1n\) is applied.
However this covering argument does incur a subexponential cost at finite \(n\) that cannot be
avoided! \fxwarning{exercise!}

\subsection{Subgaussian random variables}
\marginpar{\red{Lecture 7}}

Subgaussian random variables have tails that decay at least as fast as the
Gaussian distribution. That is
\[
    \prob{\abs{X} \ge t} \le 2 \exp\bigl(-\tfrac{t^2}{2\sigma^2}\bigr).
\]
Recall that we obtained this bound on the tails of the normal distribution using
a Chernoff bound and the CGF of the normal distribution in \eqref{ex: chernoff
bound for normal distribution}.
An important example of subgaussian random variables
are bounded random variables, since the probability 
that \(\abs{X} \ge t\) eventually drops to zero.
Unfortunately tail bounds like this do not interact
well with summation of random variables. We would therefore prefer bounds on the
CGF instead to get tight bounds on the convergence rate of the mean. We will
start with such bounds and find out later (Theorem \ref{thm: subgaussian
equivalences}) that this is actually equivalent to the tail bound above.

\begin{definition}[Variance proxy]
    Let \(X\) be a \underline{centered} random variable, then
    the \emph{variance proxy} of \(X\) is defined as
    \[
        \varProx{X} \coloneq \min\set{\lambda \ge 0 : \CGF_X(t) \le \tfrac{\lambda t^2}2 \quad \forall t}
    \]
\end{definition}

The cumulant generating function of a centered normal distribution with variance
\(\sigma^2\) is
\[
    \CGF_X(t) = \log \MGF_X(t) = \tfrac{\sigma^2t^2}{2}.
\]
The variance proxy \(\varProx{X}\) therefore coincides with the variance \(\sigma^2\) in
this case.


\begin{lemma}[Basic properties of variance proxy]
    Let \(X\) and \(Y\) be \underline{centered} random variables with variance proxies \(\varProx{X}\) and \(\varProx{Y}\), then
    \begin{enumerate}[label=(\roman*), noitemsep]
        \item\label{it: var proxy scaling} \(\varProx{a X} = a^2 \varProx{X}\) for any \(a\in \real\), in particular \(\varProx{-X} = \varProx{X}\),
        \item\label{it: variance proxy independent sums} Let \(Y\) be \underline{independent} from \(X\), then
        \[
            \varProx{X+Y} \le \varProx{X} + \varProx{Y}.
        \]
    \end{enumerate}
    In particular for the mean \(\mean{X}_n\) of independent centered \(X_i\),
    we have
    \begin{equation}
        \label{eq: variance proxy of the mean}
        \varProx{\mean{X}_n} \le \frac1{n^2} \sum_{i=1}^n \varProx{X_i}
        \overset{\substack{\text{identical}\\\text{distributed}}}= \frac{\varProx{X_1}}n.
    \end{equation}
\end{lemma}
\begin{proof}
    The first property follows from the scaling property of the CGF, that is
    \[
        \CGF_{a X}(t) = \CGF_X(a t) \le \tfrac{\varProx{X} (a t)^2}2 = \tfrac{(a^2 \varProx{X}) t^2}2.
    \]
    of course by definition of the variance proxy as a minimum we only obtain
    \(\varProx{aX} \le a^2 \varProx{X}\). To get the lower bound use the same
    argument with \(a^{-1}\) to obtain \(\varProx{X} \le a^{-2} \varProx{aX}\).
    The sum property follows from the identical
    sum property of the CGF for independent random variables, that is
    \[
        \CGF_{X+Y}(t) = \CGF_X(t) + \CGF_Y(t) \le \tfrac{\varProx{X} t^2}2 + \tfrac{\varProx{Y} t^2}2 = \tfrac{(\varProx{X} + \varProx{Y}) t^2}2.
        \qedhere
    \]

\end{proof}

The variance proxy can then be used for the following tail bound:

\begin{mdframed}[innertopmargin=0pt]
\begin{theorem}[General Hoeffding inequality]
    \label{thm: general hoeffding inequality}
    Let \(X_1, \ldots, X_n\) be independent random variables in \(\real\) with
    variance proxies \(\varProx{X_i}\), then for all \(\epsilon > 0\)
    \[
        \prob[\Big]{\mean{X}_n - \E{\mean{X}_n} \ge \epsilon}
        \le \exp\Bigl(-\frac{n \epsilon^2}{2\frac1n\sum_{i=1}^n \varProx{X_i}}\Bigr).
    \]
\end{theorem}
\end{mdframed}
\begin{proof}
    It is sufficient to prove for a centered random variable \(X\) with variance
    proxy \(\varProx{X}\) that
    \begin{equation}
        \label{eq: var prox tail bound}    
        \prob{X \ge \epsilon} \le \exp\bigl(-\tfrac{\epsilon^2}{2\varProx{X}}\bigr),
    \end{equation}
    because an application to \(X = \mean{X}_n - \E{\mean{X}_n}\) with
    \eqref{eq: variance proxy of the mean} would then give the desired result.
    For \eqref{eq: var prox tail bound} we simply use a Chernoff bound. 
    Specifically, by definition of the variance proxy \(\CGF_X(t) \le \frac{\varProx{X} t^2}2\) for all \(t\in \real\).
    We therefore obtain the following bound for the convex conjugate
    \begin{equation}
        \label{eq: bound on cgf star by variance proxy}    
        \CGF_X^*(a)
        = \sup_{t\in \real} \bigl(t a - \CGF_X(t)\bigr)
        \ge \sup_{t\in \real} \bigl(t a - \tfrac{\varProx{X} t^2}2\bigr)
        = \tfrac{a^2}{2\varProx{X}}.
    \end{equation}
    The last equation follows by the same optimization step of the quadratic
    polynomial in \(t\) as we performed when we calculated this bound for the
    normal distribution in Example \ref{ex: chernoff bound for normal
    distribution}. Then clearly by Chernoff's bound (Thm.~\ref{thm: chernoff bound in R^d})
    \[
        \prob{X\ge \epsilon} = \prob{X \in [\epsilon , \infty)}
        \overset{\text{Thm.~\ref{thm: chernoff bound in R^d}}}\le \exp\bigl(-\inf_{a \ge \epsilon}\CGF_X^*(a)\bigr)
        \overset{\text{\eqref{eq: bound on cgf star by variance proxy}}}\le \exp\bigl(-\tfrac{\epsilon^2}{2\varProx{X}}\bigr).
        \qedhere
    \]
\end{proof}

Of course the variance proxy is only useful if we can actually calculate it for
interesting random variables. The following lemma shows that for bounded random
variables the variance proxy may be bounded.

\begin{mdframed}[innertopmargin=0pt]
\begin{lemma}[Hoeffding's Lemma]
    \label{lem: hoeffding's lemma}
    Let \(X\) be a centered, bounded random variable, that is 
    \(\E{X}=0\) and \(X\in [a,b]\) with probability one. Then
    \[
        \varProx{X} \le \tfrac{(b-a)^2}4.
    \]
\end{lemma}
\end{mdframed}
\begin{proof}
    The first step of the proof is to show that the variance proxy is largest
    for the distribution that only takes the boundary values \(a\) and \(b\).
    For this we use the convexity of the exponential function. 
    Choosing \(p_x=\frac{x-b}{b-a}\in [0,1]\) for \(x\in [a,b]\) implies
    \[
        \exp(t x) \le p_x \exp(t b) + (1-p_x) \exp(t a).
    \]
    Plugging in \(X\) and taking the expectation this implies
    \[\begin{aligned}
        \E{\exp(t X)}
        &\le \E{p_X} \exp(t b) + \E{1-p_X} \exp(t a)
        \\
        &= \E{\tfrac{X-a}{b-a}} \exp(t b) + \E{\tfrac{b-X}{b-a}} \exp(t a)
        \\
        &= \tfrac{-a}{b-a}\exp(t b) + \tfrac{b}{b-a}\exp(t a)
        &&= \E{\exp(t Y)}
    \end{aligned}
    \]
    Observe that the last line is exactly the MGF of a random variable \(Y\)
    with
    \[
        \prob{Y=b} = p,
        \qquad
        \prob{Y=a} = 1-p,
        \qquad p \coloneq \tfrac{-a}{b-a}.
    \]
    Note that \(a\le 0=\E{X} \le b\) and thus \(p\in [0,1]\).
    Taking the logarithm on both sides we obtain
    \begin{equation}
        \label{eq: rademacher are worst offender}
        \CGF_X(t) \le \CGF_Y(t) \overset{!}\le \tfrac{\lambda^2 t^2}{2}
        \qquad \forall t\ge 0.
    \end{equation}
    The smallest such \(\lambda\) is the variance proxy of \(Y\).
    If we prove it is bounded by \(\frac{(b-a)^2}{4}\) we finish the proof.
    We begin by calculating the cumulant generating function of \(Y\) 
    \[\begin{aligned}
        \CGF_Y(t)
        = \log(p e^{t b} + (1-p) e^{t a})
        &= \log\bigl(e^{ta}(p e^{t (b-a)} + 1-p)\bigr)
        \\
        &= ta + \log(pe^{t(b-a)}+1-p).
    \end{aligned}
    \]
    Taylor's theorem with mean value reminder implies for some \(\xi\)
    \[
        \CGF_Y(t) = \CGF_Y(0) + \CGF_Y'(0)t + \CGF_Y''(\xi)\tfrac{t^2}{2}.
    \]
    So \(\CGF_Y(0) = \CGF_Y'(0)=0\) and \(\abs{\CGF_Y''(\xi)} \le \lambda\)
    would do the job to get \eqref{eq: rademacher are worst offender}. To
    make the calculations easier we would like to assume \((b-a)=1\).
    We can do so without loss of generality by the change of variables
    \(h=t(b-a)\), where \(p= \frac{-a}{b-a}\) implies \(-ph= ta\)
    and thus
    \[
        \CGF_Y(t)
        = -ph +  \log\bigl(p e^{h} + 1-p\bigr) \eqcolon \phi(h).
    \]
    Since \(\CGF_Y'(t) = \phi'(h) \frac{dh}{dt}\) with \(\frac{dh}{dt} =
    (b-a)\) it is sufficient to calculate the derivatives of \(\phi\) given by 
    \begin{align*}
        \phi'(h) &= \frac{pe^h}{pe^h +1-p} -p,
        \\
        \phi''(h) &= \frac{pe^h}{pe^h +1-p} - \frac{(pe^h)^2}{(pe^h +1-p)^2} = q(h)(1-q(h))
    \end{align*}
    with \(q(h) \coloneqq \frac{pe^h}{pe^h +1-p}\in [0,1]\) for all \(h\in \real\). Since \(q(1-q)\) is maximized at
    \(q=\frac12\) by the value \(\frac14\) we have \(\abs{\phi''(h)} \le \frac14\) for all \(h\in \real\).
    This implies
    \begin{alignat*}{2}
        \CGF_Y(0) &= \phi(0) &&= 0
        \\
        \CGF_Y'(0) &= \phi'(0)(b-a) &&= 0
        \\
        \abs{\CGF_Y''(\xi)} &= \abs{\phi''(\xi)}(b-a)^2 &&\le \tfrac{(b-a)^2}4
    \end{alignat*}
    In the last equation we used the chain rule again to obtain \(\CGF_Y''(\xi) = \phi''(\xi)(b-a)^2\). 
    A comparison of the the Taylor expansion with \eqref{eq: rademacher are worst
    offender} finishes the proof.
\end{proof}

\begin{mdframed}[innertopmargin=0pt]
\begin{corollary}[Hoeffding inequality for bounded random variables]
    \label{cor: hoeffding inequality for bounded random variables}
    Let \(X_1, \ldots, X_n\) be independent random variables contained in
    \([a,b]\), then we have
    \[
        \prob[\Big]{\abs[\big]{\mean{X}_n - \E{\mean{X}_n}} \ge \epsilon}
        \le 2 \exp\Bigl(-\frac{2n\epsilon^2}{(b-a)^2}\Bigr)
    \]
\end{corollary}
\end{mdframed}
\begin{proof}
    Observe that \(\varProx{X_i} \le \frac{(b-a)^2}4\) for all \(i\) by Hoeffding's Lemma and thus
    \[
        \prob[\Big]{\mean{X}_n - \E{\mean{X}_n} \ge \epsilon}
        \le \exp\Bigl(-\frac{n \epsilon^2}{2\frac1n\sum_{i=1}^n \varProx{X_i}}\Bigr)
        \le \exp\Bigl(-\frac{2n\epsilon^2}{(b-a)^2}\Bigr).
    \]
    Since \(-X_i\) is contained in \([-b, -a]\) we also have \(\varProx{-X_i} \le \frac{(b-a)^2}4\)
    and obtain the same bound for \(-(\mean{X}_n - \E{\mean{X}_n})\). This leads
    to the factor \(2\) by a union bound.
\end{proof}

Recall that we have promised that this variance proxy is actually equivalent
the tail bound
\[
    \prob{\abs{X} \ge t} \le 2 \exp\bigl(-\tfrac{t^2}{2\sigma^2}\bigr)
\]
for some \(\sigma^2\).  Observe that taking \(n=1\) in the general Hoeffding
inequality (Theorem \ref{thm: general hoeffding inequality}), we already get
this bound with \(\sigma^2 = \varProx{X}\) and consequently every distribution that satisfies a Hoeffding's
inequality must have such a tail. In the following we will see that
this tail decay and a finite variance proxy are actually equivalent among other
equivalent characterizations of subgaussian random variables.

\begin{mdframed}[innertopmargin=0pt]
\begin{theorem}[Subgaussian characterization]
    \label{thm: subgaussian equivalences}
    Let \(X\) be a random variable. The following are equivalent:
    \begin{enumerate}[label=(\roman*), noitemsep]
        \item\label{it: subgaussian tail decay}
        \textbf{Tail decay:} There exists \(K_1 > 0\) such that the tails of \(X\) satisfy
        \[
            \prob{\abs{X} \ge t} \le 2 \exp\bigl(-\tfrac{t^2}{K_1^2}\bigr)
            \qquad\text{for all } t > 0.
        \]
        \item\label{it: subgaussian moment decay}
        \textbf{Moment decay:} There exists \(K_2 > 0\) such that the moments of \(X\) satisfy
        \[
            \E{\abs{X}^p} \le K_2^p p \Gamma(\tfrac{p}2)
            \qquad\text{for all } p\ge 1.
        \]
        or equivalently,
        \[
            \norm{X}_{L^p} = \bigl(\E{\abs{X}^p}\bigr)^{1/p} \le K_2' \sqrt{p}
            \qquad\text{for all } p\ge 1.
        \]
        \item\label{it: subgaussian X^2 sub-exponential}
        \textbf{\(X^2\) sub-exponential:} There exists \(K_3 >0\) such that the MGF of \(X^2\) satisfies
        \[
            \MGF_{X^2}(t^2) = \E{\exp(t^2 X^2)} \le \exp(K_3^2 t^2)
            \qquad\text{for all } t \text{ with } \abs{t} < \tfrac1{K_3}.
        \]
        \item\label{it: subgaussian orlicz norm}
        \textbf{Orlicz norm:} There exists \(K_4 > 0\) such that
        \[
            \E{\exp\bigl(\tfrac{X^2}{K_4^2}\bigr)} \le 2.
        \]
        In other words: The Orlicz norm (Definition \ref{def: orlicz norm})
        corresponding to the convex function \(\psi_2(x) \coloneq \exp(x^2) - 1\) is finite
        \[
            \norm{X}_{\psi_2} \coloneq \inf\set{t > 0 :
            \E{\psi_2(\tfrac{\abs{X}}t)} \le 1}
            \le K_4 < \infty.
        \]
    \end{enumerate}
    If \(\E{X} = 0\), then these are also equivalent to
    \begin{enumerate}[label=(\roman*), resume, noitemsep]
        \item\label{it: subgaussian variance proxy}
        \textbf{Variance proxy:} There exists \(K_5 > 0\) such that the MGF of \(X\) satisfies
        \[
            \CGF_X(t) \le K_5^2 t^2
            \qquad\text{for all } t\in \real.
        \]
        That is \(\varProx{X} \le 2K_5^2 < \infty\).
    \end{enumerate} 
    The parameters \(K_i>0\) appearing in these statements differ from
    each other by at most a constant factor that is independent from the
    specific \(X\).
    \begin{center}
        \hspace{-1em}
        \scalebox{0.9}{
        \def\arraystretch{1.4}
        \begin{tabular}{r r | c c c c c c}
            & \multicolumn{1}{c}{} & \multicolumn{6}{c}{\normalfont have}
            \\
            & &\(K_1\) & \(K_2\) & \(K_2'\) & \(K_3\) & \(K_4\) & \(K_5\)
            \\
            \cline{2-8}
            \multirow{6}{*}{\rotatebox{90}{\normalfont want}}
            & \(K_1\)
            & \(\textcolor{gray}1\)
            & \textcolor{gray}{\(\sqrt{3}\)}
            & \textcolor{gray}{\(\frac{2\sqrt{e}}{\sqrt{\ln 2}}\)}
            & \textcolor{gray}{\(\frac1{\sqrt{\ln 2}}\)}
            & \textcolor{magenta}{\(1\)}
            & \textcolor{blue}{\(2\)}
            \\
            & \(K_2\)
            & \textcolor{magenta}{\(1\)}
            & \(\textcolor{gray}1\)
            & \textcolor{gray}{\(\frac{2\sqrt{e}}{\sqrt{\ln 2}}\)}
            & \textcolor{gray}{\(\frac1{\sqrt{\ln 2}}\)}
            & \textcolor{gray}{\(1\)} & \textcolor{gray}{\(2\)}
            \\ 
            & \(K_2'\)
            & \textcolor{gray}{\(\frac{3}{\sqrt{2}}\)}
            & \textcolor{magenta}{\(\frac{3}{\sqrt{2}}\)}
            & \(\textcolor{gray}1\)
            & \textcolor{gray}{\(\frac3{\sqrt{2\ln 2}}\)}
            & \textcolor{gray}{\(\frac{3}{\sqrt{2}}\)}
            & \textcolor{gray}{\(3\sqrt{2}\)}
            \\
            & \(K_3\)
            & \textcolor{gray}{\(\frac{6\sqrt{e}}{\sqrt{2}}\)}
            & \textcolor{gray}{\(\frac{6\sqrt{e}}{\sqrt{2}}\)}
            & \textcolor{magenta}{\(2\sqrt{e}\)}
            & \(\textcolor{gray}1\)
            & \textcolor{gray}{\(\frac{6\sqrt{e}}{\sqrt{2}}\)}
            & \textcolor{gray}{\(\frac{12\sqrt{e}}{\sqrt{2}}\)}
            \\
            & \(\mathllap{\norm{X}_{\psi_2}\! \le}K_4\)
            & \textcolor{gray}{\(\sqrt{3}\)}
            &\(\sqrt{3}\)
            & \textcolor{gray}{\(\frac{2\sqrt{e}}{\sqrt{\ln 2}}\)}
            & \textcolor{magenta}{\(\frac1{\sqrt{\ln 2}}\)}
            & \(\textcolor{gray}1\)
            & \textcolor{gray}{\(2\sqrt{3}\)}
            \\
            & \(\mathllap{\sqrt{\frac{\varProx{X}}2} \le} K_5\)
            & \textcolor{gray}{\(\frac{6\sqrt{e}}{\sqrt{2}}\)}
            & \textcolor{gray}{\(\frac{6\sqrt{e}}{\sqrt{2}}\)}
            & \textcolor{gray}{\(2\sqrt{e}\)}
            & \textcolor{blue}{\(1\)}
            & \textcolor{gray}{\(\frac{6\sqrt{e}}{\sqrt{2}}\)}
            & \(\textcolor{gray}1\)
        \end{tabular}
        \hspace{0.1em}
        \begin{minipage}{0.32\textwidth}
        If you \underline{want} \(K_i\) and \underline{have}
        \(K_j\), then the table lists a \(c_{ij}\) such that
        \[
            K_i \le c_{ij} K_j.
        \]
        The \textcolor{gray}{gray} entries follow from the other entries. The \textcolor{magenta}{magenta}
        entries prove equivalence of \(K_1\) to \(K_4\) and the
        \textcolor{blue}{blue} entries tie in \(K_5\).
        \end{minipage}
        }
    \end{center}
\end{theorem}
\end{mdframed}

\begin{proof}[Proof {\citep[Prop.~2.5.2]{vershyninHighDimensionalProbabilityIntroduction2018}}]
    The equivalence clearly follows if we bound the \(c_{ij}\) for sufficient
    combinations of \(i\) and \(j\) such that loops may be formed.

    ``\ref{it: subgaussian tail decay} \(\Rightarrow\) \ref{it: subgaussian moment decay}''
    Let us start with \(K_1\)
    \[\begin{aligned}
        \E{\abs{X}^p}
        &= \int_0^\infty \prob{\abs{X}^p \ge u} du
        \mathrlap{= \int_0^\infty \prob{\abs{X} \ge u^{1/p}} du}
        \\
        &= \int_0^\infty \prob{\abs{X} \ge t} p t^{p-1} dt
        & (u=t^p)
        \\
        &\le p \int_0^\infty 2\exp(-\tfrac{t^2}{K_1^2}) t^{p-1} dt
        \\
        &= K_1^p p \int_0^\infty e^{-s} s^{\frac{p}2 - 1} ds
        & \Bigl(s = (\tfrac{t}{K_1})^2, dt = \tfrac{K_1}{2} s^{-\frac12} ds\Bigr)
        \\
        \overset{\text{def.}}&= K_1^p p \Gamma(\tfrac{p}2).
    \end{aligned}
    \]
    So clearly \(K_2 \le K_1\) (i.e. \(c_{21} = 1\)). Observe that
    we could have assumed \(K_1=1\) without loss of generality by bounding the
    scaled random variable \(X/K_1\). We will do so from now on.
    
    To get \(K_2'\) from \(K_2\) we use a bound on the Gamma function,
    \(\Gamma(x)\le 3x^x\) for
    \(x\ge \tfrac12\) \citep{vershyninHighDimensionalProbabilityIntroduction2018}. Then we have
    \[
        \bigl(p\Gamma(\tfrac{p}2)\bigr)^{\frac1p}
        \le \bigl(3p(\tfrac{p}2)^{\frac{p}2}\bigr)^{\frac1p}
        = \sqrt{\frac{p}2} (3p)^{\frac1p}
        \le \frac{\max_{p\ge 1} f(p)}{\sqrt{2}} \sqrt{p}
    \]
    with \(f(x) = (3x)^{\frac1x} = \exp(\frac1x\ln(3x))\). Its derivative is given by
    \[
        f'(x) = f(x) \tfrac1{x^2}(1- \ln(3x)) \le 0 \qquad \text{for } x \ge 1.
    \]
    Consequently the maximizer is \(f(1) = 3\) and thus
    \(K_2' \le K_2 \frac{3}{\sqrt{2}}\).

    ``\ref{it: subgaussian moment decay} \(\Rightarrow\) \ref{it: subgaussian X^2 sub-exponential}''
    To get \(K_3\) from \(K_2'\) observe that by the series expansion of the exponential
    function
    \[
        \E{\exp(t^2 X^2)}
        = \E[\Big]{1+\sum_{k=1}^\infty \frac{(t^2 X^2)^k}{k!}}
        =1+ \sum_{k=1}^\infty \frac{t^{2k} \E{X^{2k}}}{k!}
    \]
    using the assumption \(\E{X^{2k}} \le (2k)^{\frac{2k}2}\) with \(K_2'=1\) and
    Stirling's approximation \(k! \ge (\tfrac{k}e)^k\) this results in
    \[
        \E{\exp(t^2 X^2)}
        \le 1+ \sum_{k=1}^\infty \frac{t^{2k} (2k)^k}{(k/e)^k}
        = 1 + \sum_{k=1}^\infty (2e t^2)^k
        = \frac1{1-2e t^2},
    \]
    assuming \(2e t^2 < 1\). Further assuming \(2e t^2 \le \frac12\) and using
    the inequality \(\frac1{1-x} \le e^{2x}\) for \(x\in [0, \frac12]\) we have
    \[
        \E{\exp(t^2 X^2)} \le \exp(4e t^2) \qquad \text{for } t < \tfrac1{2\sqrt{e}}.
    \]
    Consequently \(K_3 \le 2\sqrt{e} K_2'\).

    ``\ref{it: subgaussian X^2 sub-exponential} \(\Rightarrow\) \ref{it: subgaussian orlicz norm}''
    Let \(c=K_3/\sqrt{\ln 2}\) such that \(\frac1c \le \frac1{K_3}\)
    and therefore
    \[
        \E{\exp{X^2/c^2}} \le \exp(K_3^2/c^2) = \exp(\ln 2) = 2.
    \]
    Consequently \(K_4 \le c= \frac{K_3}{\sqrt{\ln 2}}\).
    
    ``\ref{it: subgaussian orlicz norm} \(\Rightarrow\) \ref{it: subgaussian
    tail decay}'' Using Markov's inequality (Theorem \ref{thm: markov
    inequality}) on \(h(x) = e^{x^2}\) we have for \(t > 0\)
    \[
        \prob{\abs{X} \ge t} \le \tfrac{\E{\exp(X^2)}}{\exp(t^2)}
        \le 2 \exp(-t^2) \qquad \text{assuming } K_4 = 1.
    \]
    Consequently, \(K_1 \le K_4\).

    ``\ref{it: subgaussian X^2 sub-exponential}'' \(\Rightarrow\) ``\ref{it: subgaussian variance proxy}''
    We assume \(K_3=1\) again
    and use the inequality \(e^x \le x + e^{x^2}\) for \(x\in \real\) to get
    \[
        \E{\exp(t X)}
        \le \E{t X + \exp(t^2 X^2)}
        \overset{\E{X}=0}= \E{\exp(t^2 X^2)}
        \le \exp(t^2).
    \]
    The last inequality only holds for \(\abs{t} < \frac1{K_3} = 1\). Now assume
    \(\abs{t} \ge \frac1{K_3}\). Using the inequality \(2tx \le t^2 + x^2\), we get
    \[
        \E{\exp(t X)}
        \le e^{\frac{t^2}2}\E{e^{\tfrac{X^2}2}}
        \le e^{\frac{t^2}2} e^{\tfrac12} \overset{\abs{t}\ge 1}\le e^{t^2}
    \]
    The penultimate inequality follows from \(\frac12 \le \frac1{K_3} =1\).
    As we have covered both case we get \(K_5 \le K_3\).

    ``\ref{it: subgaussian variance proxy} \(\Rightarrow\) \ref{it: subgaussian tail decay}''
    Since \(\varProx{X} \le 2K_5^2\) we have by the general Hoeffding inequality for \(n=1\), that is \eqref{eq: var prox tail bound}
    \[
        \prob{X \ge t} \le \exp\bigl(-\tfrac{t^2}{4K_5^2}\bigr).
    \]
    Taking the same bound for \(-X\) we get \(K_1^2 \le 4K_5^2\) and thus \(K_1 \le 2 K_5\).
\end{proof}

\begin{remark}[Sums of (in)dependent subgaussian random variables]
    Since the Orcliz norm is a norm
    (Theorem \ref{thm: orlicz seminorm}), it holds for all (possibly dependent)
    random variables \(X\) and \(Y\)
    \[
        \norm{X+Y}_{\psi_2} \le \norm{X}_{\psi_2} + \norm{Y}_{\psi_2}.
    \]
    In comparison the sum stability of the variance proxy only holds for
    independent random variables, that is
    \[
        \varProx{X+Y} \le \varProx{X} + \varProx{Y}
        \qquad \text{for } X, Y \text{ independent}.
    \]
    This may look like an advantage of the Orcliz norm but note that \(\varProx{X}\)
    is something like a squared norm, as \(\varProx{aX} = a^2 \varProx{X}\),
    while \(\norm{aX}_{\psi_2} = \abs{a} \norm{X}_{\psi_2}\).
    Put another way, since \(\norm{X}_{\psi_2}\) is the smallest \(K_4\),
    and \(\sqrt{\varProx{X}}\) is the smallest \(K_5\) up to the constant
    \(\sqrt{2}\) due to \(\sqrt{\varProx{X}}\le 2K_5^2\), we have by
    the equivalence of constants
    \[
        \norm{X}_{\psi_2} \le c_0 \sqrt{\varProx{X}} \le c_1 \norm{X}_{\psi_2}
    \]
    and thus for independent \(X\) and \(Y\)
    \[
        \norm{X + Y}_{\psi_2}^2 \le c_0^2 \varProx{X+Y} \le c_0^2 (\varProx{X} + \varProx{Y})
        \le c_1^2 (\norm{X}_{\psi_2}^2 + \norm{Y}_{\psi_2}^2).
    \]
    This bound is generally better than if we applied the triangle inequality inside
    the square due to the removal of the cross term
    \(2\norm{X}_{\psi_2}\norm{Y}_{\psi_2}\), that is
    \[
        \norm{X}_{\psi_2}^2 + \norm{Y}_{\psi_2}^2 \le (\norm{X}_{\psi_2} + \norm{Y}_{\psi_2})^2.
    \]
    In other words: the sum stability of the squared norm is better than the sum
    stability of the norm itself. This is particularly visible for the mean in
    the case of iid \(X_i\)
    \[
        \norm{\mean{X}_n}_{\psi_2} \le c_0 \sqrt{\varProx{\mean{X}_n}}
        = c_0 \sqrt{\red{\frac1{n^2}} \sum_{i=1}^n \varProx{X_i}}
        =  \frac{c_0}{\sqrt{n}} \underbrace{\sqrt{\frac1n \sum_{i=1}^n \varProx{X_i}}}_{= \varProx{X_i}}
        \in \bigO(\tfrac1{\sqrt{n}}),
    \]
    which is a much better bound asymptotically than the general triangle inequality
    \[
        \norm{\mean{X}_n}_{\psi_2} \le \frac1n \sum_{i=1}^n \norm{X_i}_{\psi_2} = \norm{X_i}_{\psi_2} \in \bigO(1).
    \]
    \textbf{Rule of thumb}: work with the variance proxy if the variables are independent
    and with the Orlicz norm if they are not. The translation between
    these view points may however introduce constants, which can in certain
    circumstances be larger than the gains from this change of view.
\end{remark}


\subsection{Subexponential random variables}

\begin{theorem}[Subexponential characterization]
    
\end{theorem}

\begin{theorem}[Bernstein inequality]
    
\end{theorem}


\subsection*{Exercises}

\begin{exercise}[Properties of the cumulant generating function]
    \label{ex: properties cumulant generating function}
    Let \(X\) be a random variable. Then the cumulant generating function
    \(\CGF_X(t) = \log \MGF_X(t)\) has the following properties:
    \begin{enumerate}[label=(\roman*), noitemsep]
        \item\label{it: cgf 0=0} \(\CGF_X(0) = 0\).
        \item\label{it: cgf linear} \(\CGF_{aX + b}(t) = \scp{b, t} + \CGF_X(\scp{a, t})\) for any \(a \in \real\), \(b,t\in \hilbert\).
        \item\label{it: cumulant independent sums} Let \(Y\) be \underline{independent} from \(X\), then
        \[
            \CGF_{X+Y}(t) = \CGF_X(t) + \CGF_Y(t)
        \]
    \end{enumerate}
    If the moment generating function is finite in a neighborhood of zero,
    then \(\CGF_X\) is infinitely often differentiable and
    \begin{enumerate}[label=(\roman*), resume, noitemsep]
        \item \(\CGF_X'(0) = \E{X}\),
        \item \(\CGF_X''(0) = \var{X}\)
    \end{enumerate}

    \begin{solution}
        For \ref{it: cgf 0=0}-\ref{it: cumulant independent sums} follow from the
        corresponding properties of the MGF (Exercise \ref{ex: properties moment
        generating functions}) by application of the logarithm.  For example,
        \(\CGF_X(0)=0\) follows directly from \(\MGF_X(0)=1\). For convexity
        let \(t, s\in \hilbert\) and \(\lambda \in [0,1]\). 

        The derivatives are straightforward calculations using the derivatives of the MGF
        and the chain rule due to \(\CGF_X(t) = \log \MGF_X(t)\). That is
        \begin{alignat*}{3}
            \CGF_X'(t)
            &= \frac{\MGF_X'(t)}{\MGF_X(t)}
            \qquad &\implies &\CGF_X'(0) = \E{X},
            \\
            \CGF_X''(t)
            &= \frac{\MGF_X''(t)}{\MGF_X(t)} - \frac{(\MGF_X'(t))^2}{(\MGF_X(t))^2}
            \qquad &\implies &\CGF_X''(0)
            \begin{aligned}[t]
            &= \E{X^2} - (\E{X})^2
            \\
            &= \var{X}.
            \end{aligned}
        \end{alignat*}
        This proves the claims.
    \end{solution}
\end{exercise}

\begin{definition}[Cumulants]
    If they exist, the ``cumulants'' \(\cumulant_n^X\) are the coefficients of the Taylor
    series of the cumulant generating function
    \[
        \CGF_X(t)
        = \sum_{n=1}^\infty \cumulant_n^X \frac{t^n}{n!}
        = \sum_{n=1}^\infty \CGF_X^{(n)}(0) \frac{t^n}{n!}.
    \]
\end{definition}

Due to property \ref{it: cumulant independent sums} in Exercise \ref{ex:
properties cumulant generating function} of the cumulant generating function,
cumulants have the nice property that they ``accumulate'' when indpendent
random variables are added, that is for \(X\) and \(Y\) independent
\[
    \cumulant_n^{X+Y} = \cumulant_n^X + \cumulant_n^Y.
\]
This property is the origin of their name and what makes them ``nicer'' compared
to the moments that make up the coefficients of the moment generating function.
Bienaymé (Lemma \ref{lem: Bienaymé}) represents this property for the second
order cumulants, since \(\cumulant_2^X = \var{X}\).

\begin{remark}[Cumulants of normal distribution]
    Recall that the cumulant generating function of the 
    normal distribution is a quadratic polynomial
    \[
        \CGF_X(t) = \mu t + \frac12 \sigma^2 t^2.
    \]
    Thus all cumulants of order three and higher are zero! It is the only
    distribution with this property. One can show that the cumulants of order
    three and higher vanish for \(\sqrt{n}\, \mean{X}_n\) as \(n\to\infty\)
    (Exercise \ref{ex: cumulant CLT}) and thereby proof a weak form of the
    central limit theorem.
\end{remark}
\begin{exercise}[Cumulant CLT]
    \label{ex: cumulant CLT}
    Let \(X_1, X_2, \ldots\) be independent copies of a random variable
    \(X\) with mean \(\mu\), variance \(\sigma^2\) and finite moments of all orders.
    Let \(\mean{X}_n = \frac1n \sum_{i=1}^n X_i\) be the empirical mean.
    Show that the cumulants \(\cumulant_k^{\sqrt{n}(\mean{X}_n - \mu)}\)
    of the random variable \(\sqrt{n}(\mean{X}_n - \mu)\) satisfy
    \[
        \cumulant_1^{\sqrt{n}(\mean{X}_n - \mu)} = 0,
        \qquad
        \cumulant_2^{\sqrt{n}(\mean{X}_n - \mu)} = \sigma^2,
        \qquad
        \cumulant_k^{\sqrt{n}(\mean{X}_n - \mu)} \to 0
        \quad\text{for } k \ge 3.
    \]
    Conclude that the cumulant generating function of
    \(\sqrt{n}(\mean{X}_n - \mu)\) converges to the cumulant generating function
    of the normal distribution \(\normal{0, \sigma^2}\).
\end{exercise}

\begin{exercise}[Convex conjugate]
    \label{ex: convex conjugate}
    
\end{exercise}

\paragraph*{Randomized algorithms} The theme of the
next two exercises is that algorithms can sometimes benefit from randomization.

\begin{exercise}[Monte Carlo \Coffeecup]
    \label{ex: monte carlo}
    In this exercise we will devise a way to compute the integral
    \[
        I \coloneq \int_{[0,1]^\dims} f(x) dx
    \]
    of a continuous function \(f\colon [0,1]^\dims \to \real\).  Let
    \((U_n)_{n\in \nat}\) be independent and uniformly distributed on
    \([0,1]^\dims\).
    \begin{enumerate}[label=(\alph*), noitemsep]
        \item Prove that \(\hat I_n \coloneq \frac1n\sum_{i=1}^n f(U_i) \overset\as\to I\).
        \begin{solution}
            Since \(I= \E{f(U)}\) this is the strong law of large numbers,
        \end{solution}

        \item Using \(\norm{f}_\infty = \sup_{x\in [0,1]^\dims} \abs{f(x)}\) Prove that
        \[
            \prob[\big]{\abs{\hat I_n - I} > \epsilon} \le 2\exp\Bigl(-\frac{n \epsilon^2}{2\norm{f}_\infty^2}\Bigr)
        \]
        \begin{solution}
            Since \(f\) is continuous, it is bounded on the compact set \([0,1]^\dims\)
            and thereby \(f(U) \in [-\norm{f}_\infty, \norm{f}_\infty]\).
            By Hoeffding's inequality (Corollary \ref{cor: hoeffding inequality for bounded random variables})
            for bounded random variables we therefore have the claim.
        \end{solution}

        \item Let \(n = n(\epsilon, \delta)\) be the smallest \(n\) such that the following
        PAC bound holds
        \[
            \prob[\big]{\abs{\hat I_n - I} > \epsilon} \le \delta
        \]
        Prove that
        \[
            n(\epsilon, \delta)
            \le \ceil[\Big]{2\norm{f}_\infty^2\frac{\ln(\tfrac2\delta)}{\epsilon^2}}
            \in \bigO\Bigl(\norm{f}_\infty^2\frac{\log(\frac1\delta)}{\epsilon^2}\Bigr)
        \]
        \begin{remark}[Accuracy vs confidence]
            \label{rem: accuracy vs confidence}
            The accuracy \(\epsilon\) is expensive and grows like
            \(\frac1{\epsilon^2}\) while the confidence \(1-\delta\) is cheap and grows only
            like \(\log(\frac1\delta)\). Note that deterministic numerical quadrature
            formulas suffer from the curse of dimensionality causing a numerical
            complexity of order \(\epsilon^{-\frac{\dims}k}\) \citep[Sec.~1.3.12, Prop.~1]{novakDeterministicStochasticError1988}, where \(k\) is the
            number of continuous derivatives of \(f\). In comparison, the 
            computational complexity of the Monte Carlo method is independent of
            the dimension \(\dims\). This means that Monte Carlo methods become
            interesting for high-dimensional integration problems -- specifically,
            when \(\frac{\dims}k \ge 2\).
        \end{remark}
        \begin{solution}
            By the previous exercise it is sufficient to have
            \begin{align*}
                2\exp\Bigl(-\frac{n\epsilon^2}{2\norm{f}_\infty^2}\Bigr) \le \delta
                &\iff -\frac{n\epsilon^2}{2\norm{f}_\infty^2} \le \ln(\tfrac{\delta}2)
                \\
                &\iff \ln(\tfrac{2}\delta) \le \frac{n \epsilon^2}{2\norm{f}_\infty^2}
                \\
                &\iff 2\norm{f}_\infty^2 \frac{\ln(\tfrac2\delta)}{\epsilon^2} \le n
            \end{align*}
            Since this is a sufficient condition, the smallest \(n(\delta, \epsilon)\)
            is smaller than the first natural number that satisfies it, leading
            to the claim.
        \end{solution}
    \end{enumerate}
\end{exercise}


In the following exercise we will show how to speed-up the computation of the
scalar product in high dimension using a randomized approach.
Since matrix multiplication is made up of many scalar products it should be
clear that this is only a tiny glimps into the field of randomized
linear algebra \citep[e.g.][]{murrayRandomizedNumericalLinear2023}.

\begin{exercise}[Randomized linear algebra \Coffeecup]
    Let \(x,y \in \real^\dims\), where the dimension \(\dims\) is really large.
    The computation of \(\scp{x,y}\) clearly costs \(\dims\) multiplications and
    additions. We will try to speed it up using a random sample.
    Specifically, let \(Z_1, \ldots Z_n\) be independently drawn from
    \(\frac1\dims\sum_{i=1}^\dims \delta_{x_iy_i}\), that is
    \(\prob{Z_1 = x_iy_i} = \frac1\dims\) for \(i=1, \ldots, \dims\).
    \begin{enumerate}[label=(\alph*), noitemsep]
        \item Prove that \(\dims \E{Z_1} = \scp{x,y}\),
        \begin{solution}
            By definition we have
            \[
                \E{Z_1} = \frac1\dims\sum_{i=1}^\dims  x_i y_i = \frac1\dims \scp{x,y}.
                \qedhere
            \]            
        \end{solution}

        \item Assuming \(\scp{x,y} \neq 0\), prove the following estimator
        \[
            \hat S_n \coloneq \frac\dims{n} \sum_{i=1}^n Z_i,
        \]
        converges against the scalar product in relative terms, that is
        \[
            \frac{\hat S_n}{\scp{x,y}} \overset\as\to 1 \qquad (n \to \infty).
        \]
        \begin{solution}
            We have \(\frac{\hat S_n}{\scp{x,y}} = \frac1n \sum_{i=1}^n Z_i /
            \E{Z_1}\) and thus the claim by the strong law of large numbers.
        \end{solution}

        \item For \(\hat S_n\) to be an effective speedup of the
        computation of the scalar product it is necessary that \(n\le \dims\). So
        clearly \(n\to \infty\) is not an interesting statement.
        Let \(m = \max_i \abs{x_i y_i}\) and show that
        \[
            \prob[\Big]{\abs[\Big]{\frac{\hat S_n}{\scp{x,y}} - 1} > \epsilon}
            \le 2\exp\Bigl(-n \epsilon^2 \frac{\E{Z_1}^2}{8 m^2}\Bigr).
        \]
        \textbf{Hint:} Corollary \ref{cor: hoeffding inequality for bounded random variables}.
        \begin{solution}
            Since \(Z_i \in [-m,m]\) for \(i=1, \ldots, n\) we can apply
            Hoeffding's inequality for bounded random variables (Corollary \ref{cor:
            hoeffding inequality for bounded random variables}) to get
            \begin{align*}
                \prob[\Big]{\abs[\Big]{\frac{\hat S_n}{\scp{x,y}} - 1} > \epsilon}
                &=\prob[\Big]{\abs[\Big]{\frac1n\sum_{i=1}^n Z_i - \E{Z_1}} > \epsilon \abs{\E{Z_1}}}
                \\
                &\le 2\exp\Bigl(-n \epsilon^2 \frac{\E{Z_1}^2}{2(2m)^2}\Bigr).
                \qedhere
            \end{align*}
        \end{solution}
        \begin{remark}[Many significant values]
            If most \(x_iy_i\) are very small except for a few \(i=\set{1,\dots, \dims}\),
            then the inner product \(\scp{x,y}\) is dominated by these values.
            This means that a random sample can be very far off if it does not
            contain these significant indices. This case is captured
            by the ratio \(\frac{\E{Z_1}}{m}\): If the average product
            \(\E{Z_1}\) is much smaller, than largest product \(m\), the sum is
            dominated by a few large entries and convergence is slow.
        \end{remark}

        \item Assume \(\frac{\E{Z_1}}{m} \ge 0.1\), that a precision of \(\epsilon
        = 10^{-4}\) is required and a \(99\%\) confidence. How big must
        \(\dims\) be such that this approach reduces the computational cost?
        How big to get a \(99.999\%\) confidence?
        \begin{solution}
            We need \(n\) to satisfy
            \[
                2\exp\Bigl(-n \epsilon^2 \frac{\E{X_1 Y_1}^2}{8 m^2}\Bigr) \le \delta 
                \iff n \ge \frac{8\ln(\frac1{2\delta})}{\epsilon^2 \frac{\E{X_1 Y_1}^2}{m^2}}.
            \]
            Since \(\E{X_1 Y_1}^2 \ge 0.1\) we have that
            \[
                n \ge \frac{80\ln(\frac1{2\delta})}{\epsilon^2}
            \]
            is sufficient to get the required confidence. We now simply
            plug in the numbers. For \(\delta = 0.01 = 1\%\) we have
            \[
                n \ge \frac{80\ln(\frac1{0.02})}{10^{-8}} \approx 3.12 \cdot 10^{10}.
            \]
            For \(\delta = 0.00001 = 0.001\%\) we have
            \[
                n \ge \frac{80\ln(\frac1{0.00002})}{10^{-8}} \approx 8.66 \cdot 10^{10}.
            \]
            If \(n\le \dims\) then this approach reduces the computational cost.
            So for \(\dims \ge 3.12 \cdot 10^{10}\) (or \(\dims \ge 8.66 \cdot 10^{10}\) respectively)
            this approach reduces the computational cost. As in Remark \ref{rem:
            accuracy vs confidence} the confidence has almost no effect on the
            required dimension. In contrast the precision \(\epsilon\)
            dominates the required dimension.
        \end{solution}
    \end{enumerate}
\end{exercise}

\begin{exercise}[Single issue polling \Coffeecup]
    \label{ex: single issue polling} 
    Let \(x_1,\ldots, x_N \in \set{\text{no}, \text{yes}} \simeq \set{0,1}\) be the
    opinions of \(N\) people on a certain decision. In the following we will
    assume \(x_i \in \set{0,1}\). Let
    \[
        p = \frac1N\sum_{i=1}^N x_i
    \]
    be the share of people that agree with the decision.

    Assume that you can randomly sample \(n\) people from the population and
    and that they truthfully tell you their opinions \(X_1, \dots, X_n\).  That
    is \(\prob{X_i = x_j} = \frac1N\) for \(j=1, \ldots, N\) and the \(X_i\) are
    independent.
    \begin{enumerate}[label=(\alph*), noitemsep]
        \item Prove the \(X_i\) are Bernoulli distributed, specifically \(X_i \sim \bernoulli(p)\). 
        \begin{solution}
            We have
            \begin{align*}
                \prob{X_i = 1}
                &= \sum_{j=1}^N \prob{X_i = 1, X_i = x_j}
                \\
                &= \sum_{j=1}^N \prob{X_i = 1 | X_i = x_j} \prob{X_i = x_j}
                = \sum_{j=1}^N x_j \frac1N = p.
                \qedhere
            \end{align*}
        \end{solution}
        \item 
        Prove that
        \[
            \prob[\Big]{\abs[\Big]{\frac1n \sum_{i=1}^n X_i - p} > \epsilon}
            \le 2\exp(-2n \epsilon^2).
        \]
        \begin{solution}
            Since \(X_i \in \set{0,1}\) we can apply Hoeffding's inequality for
            bounded random variables (Corollary \ref{cor: hoeffding inequality for
            bounded random variables}) to get the claim.
        \end{solution}

        \item Let \(\abs{p - 0.5} \ge 0.1\), how many people \(n\) do you need to sample to
        be \(99\%\) confident on the majority opinion? How many, if the opinion is more
        controversial and \(\abs{p - 0.5} \ge 0.01\)?

        \textbf{Hint:} Assume \(p \le 0.4\) without loss of generality.
        \begin{solution}
            If \(p \le 0.5 - \epsilon\) for \(\epsilon >0\) the majority opinion
            is ``no'' and we need to be \(1-\delta\) confident that \(\frac1n
            \mean{X}_n \coloneq \sum_{i=1}^n X_i < 0.5\). This requires
            \begin{align*}
                \delta \overset{!}\ge \prob[\Big]{\mean{X}_n \ge 0.5}
                &= \prob[\Big]{\mean{X}_n - p \ge 0.5 - p}
                \\
                &\le \prob[\Big]{\mean{X}_n - p \ge \epsilon}
                \\
                &\le \exp(-2n \epsilon^2)
            \end{align*}
            In the last equation we used the General Hoeffding's inequality
            Theorem \ref{thm: general hoeffding inequality} and Hoeffding's
            lemma (Lemma \ref{lem: hoeffding's lemma}) applied to \(X_i \in
            \set{0,1}\).
            Consequently, \(\exp(-2n \epsilon^2) \le \delta\) is sufficient, and therefore
            \[
                n \ge \frac{\ln(\frac1\delta)}{2\epsilon^2}
            \]
            is sufficient.  Similar arguments apply if \(p \ge 0.5 + \epsilon\).
            In our case we
            have \(\delta = 1\% = 0.01\) and \(\epsilon = 0.1\) or \(\epsilon =
            0.01\) respectively, which implies
            \[
                n \ge \begin{cases}
                    \ceil{\frac{\ln(100)}{2 \cdot 0.1^2}} = 231 & \text{if } \epsilon = 0.1 \iff \abs{p-0.5} \ge 0.1,
                    \\
                    \ceil{\frac{\ln(100)}{2 \cdot 0.01^2}} = 23026 & \text{if } \epsilon = 0.01 \iff \abs{p-0.5} \ge 0.01.
                \end{cases}
            \]
            is sufficient for \(99\%\) confidence.
        \end{solution}
    \end{enumerate}
\end{exercise}